% SCP10, Observation 6.6, source lines 1888--1909.
% Original-tensor normalized corollary for untwisted coarse periods at least three.

\begin{theorem}[Transport between physical realizations of one symmetry]%
    \label{thm:peps_g_isometric_physical_transport}
    \lean{TNLean.PEPS.IsGIsometric.exists_physicalSupportTransport,
        TNLean.PEPS.exists_graphGIsometricPhysicalSupportTransport}
    \leanok
    \uses{thm:peps_g_isometric_support_coordinates,
        thm:peps_graph_physical_support_transport}
    Let $A,B$ be G-isometric physical maps for the same finite-dimensional
    unitary representation of a finite group. There are $c_A,c_B>0$ and a
    physical map $F$ such that $FA=\sqrt{c_A/c_B}\,B$, and $F$ preserves
    inner products on $\operatorname{ran}A$. For finite graph tensors with
    this property at every vertex, the product of the local maps restricts
    to an isometry $I$ on the actual product physical range, with
    \begin{align}
        I\Psi_A=\left(\prod_v\sqrt{c_{A,v}/c_{B,v}}\right)\Psi_B.
    \end{align}
    The graph states are their actual bond contractions. All local physical
    alphabets may vary with the vertex.
\end{theorem}

\begin{proof}\leanok
    Let $J_A=c_A^{-1/2}A^*$ be the normalized adjoint and set
    $F=c_B^{-1/2}BJ_A$. On the range of $A$, the vector $J_Ax$ is invariant,
    so G-isometry of $B$ and isometry of $J_A$ imply
    $\langle Fx,Fy\rangle=\langle x,y\rangle$. Moreover,
    $FA=\sqrt{c_A/c_B}\,BP_\rho=\sqrt{c_A/c_B}\,B$.
    The local column Gram equalities imply the same inner-product equality
    for the product physical map on the product range. Expanding its action
    on the graph contraction gives the displayed state equality.
\end{proof}

\begin{definition}[Physical tensor product of coefficient vectors]%
    \label{def:peps_physical_state_product}
    \lean{TNLean.PEPS.physicalStateProduct}
    \leanok
    For physical coefficient vectors $\psi\in\mathbb C^P$ and
    $\omega\in\mathbb C^Q$, their tensor product has complete pair-indexed
    coefficients $(\psi\otimes\omega)(p,q)=\psi(p)\omega(q)$.
\end{definition}

\begin{theorem}[Normalization of a physical tensor-product identity]%
    \label{thm:peps_physical_state_normalization}
    \lean{TNLean.PEPS.norm_physicalStateProduct,
        TNLean.PEPS.norm_normalized_physicalState,
        TNLean.PEPS.LinearIsometry.normalized_physicalStateProduct}
    \leanok
    \uses{def:peps_physical_state_product}
    The physical tensor product satisfies
    $\|\psi\otimes\omega\|=\|\psi\|\|\omega\|$.
    Every nonzero vector $\psi$ has normalized vector
    $\widehat\psi=\|\psi\|^{-1}\psi$ of unit norm.
    If an isometry satisfies $Ix=r\,\psi\otimes\omega$ with $r>0$ and
    $\|\omega\|=1$, then $I\widehat x=\widehat\psi\otimes\omega$.
\end{theorem}

\begin{proof}\leanok
    The squared norm of the pair-indexed coefficients is the product of
    the two finite sums of squared absolute values. Thus
    $\|x\|=\|Ix\|=r\|\psi\|$, and the positive scalar cancels in the
    normalized equality.
\end{proof}

\begin{definition}[Original coarse physical registers and scale]%
    \label{def:peps_two_by_two_original_output}
    \lean{TNLean.PEPS.twoByTwoOriginalOutputRegrouping,
        TNLean.PEPS.twoByTwoOriginalScale}
    \leanok
    Let $\Gamma$ be the coarse torus graph. At each coarse vertex, retain
    the original physical label $s_v$ and every relative surplus half-edge
    label $r_{v,e}$. Their full physical space is identified unitarily with
    the tensor product of the original coarse physical space and the surplus
    space. For positive site factors $c_{A,v},c_{B,v}$ put
    \begin{align}
        R=\left(\prod_v\frac{\sqrt{c_{A,v}}}{\sqrt{c_{B,v}}}\right)
        (\sqrt{|G|})^{|E(\Gamma)|}.
    \end{align}
\end{definition}

\begin{theorem}[Positivity and Bell normalization]%
    \label{thm:peps_two_by_two_original_scale}
    \lean{TNLean.PEPS.norm_regularGraphNormalizedBell,
        TNLean.PEPS.twoByTwoOriginalScale_pos}
    \leanok
    \uses{def:peps_two_by_two_original_output}
    The scale $R$ is strictly positive. The residual graph Bell product
    $\Omega$ has Hilbert norm one.
\end{theorem}

\begin{proof}\leanok
    Every square root and every denominator is positive. Each surplus
    equality vector has its inverse-square-root dimension normalization,
    so their product has squared norm one and hence norm one.
\end{proof}

\begin{theorem}[Normalized renormalization with the original tensor]%
    \label{thm:peps_two_by_two_original_normalized}
    \lean{TNLean.PEPS.exists_regularTwoByTwoOriginalTensorIsometry}
    \leanok
    \uses{thm:peps_g_isometric_physical_transport,
        thm:peps_two_by_two_bundled_isometry,
        thm:peps_two_by_two_fine_support,
        thm:peps_physical_state_normalization,
        thm:peps_two_by_two_original_scale,
        thm:peps_regular_graph_surplus_isometry,
        thm:peps_regular_graph_surplus_factorization,
        thm:peps_two_by_two_nonzero}
    Let $A$ be a regular G-isometric four-leg tensor for any finite group
    and any finite physical alphabet. Suppose the two coarse periods are
    at least three and every torus bond is untwisted. Let $\Psi_f$ be its
    actual state on the doubled fine torus and $\Psi_c$ its actual state
    with the same original tensor on the coarse torus. Both states are
    nonzero. There are positive block and target site factors, local
    physical maps $F_v$, and an isometry $I$ on the derived fine physical
    support such that
    \begin{align}
        I\Psi_f         & =R\,\Psi_c\otimes\Omega,      \\
        I\widehat\Psi_f & =\widehat\Psi_c\otimes\Omega.
    \end{align}
    The complete physical operation consists of four-site physical
    grouping, the product of the maps $F_v$, and regrouping of the original
    and surplus output registers. The Bell state has norm one, and no
    ambient physical surjectivity is assumed. This is the untwisted
    periodic specialization of
    \cite[Observation~6.6]{Schuch2010PEPS}, with the original coarse tensor.
\end{theorem}

\begin{proof}\leanok
    The actual four-site block tensors are G-isometric from fine-site
    isometry. The target is the original tensor with its surplus identity
    registers, which is G-isometric for the same bundled representation.
    Physical support transport gives the first equality, with one ratio
    of local positive factors per block. The target graph contraction
    factors into the original coarse contraction and surplus equality
    vectors. Normalizing the latter contributes the Bell factor in $R$.
    The identity closure sector proves nonvanishing of both original
    states. Their exact isometric norm relation then gives the second
    equality.
\end{proof}

The physical fixed-point statement uses untwisted coarse tori with both
periods at least three. The bond-indexed geometric regrouping has a separate
all-positive-period statement; it does not remove the physical theorem's
period or closure assumptions. The exact scope is recorded in
\cite{gap:scp10_two_by_two_regular_fixed_point}.
